% Einheit 4 von 4 – Winkel von Geraden gegen Ebenen und Bogenmaße (bank/skalarprodukt-und-winkel/e4.jsonl, zusammenbau v0.1)
%% TODO Zeitmarke Einheit 4: Marken-Zeile nicht lesbar
%% TODO Zweigzeile Teil 1: was hier gelernt wird (Satz in Schülersprache aus dem Einheitstitel) – Einheit 4
\einheitenkopf[e4]{Einheit 4 von 4 · Winkel von Geraden gegen Ebenen und Bogenmaße}
\zweigzeile{Abitur GK}
\setcounter{aufgabe}{15}

%% TODO Titel: Ich-kann-Satz fehlt in der Bank (Platzhalter: Kettenname) – Nr. 16
%% TODO Anweisung: ein Satz über der ersten Teilaufgabe fehlt in der Bank – Nr. 16
\begin{aufgabe}{Winkel gegen Ebenen und Bogenmaße}
%% TODO \rechenplatz{n}: Schrittzahl des Grundfalls fehlt in der Bank (nur bei fester Schreibform, 2.3 b)
\begin{teile}
\teil Der Winkel zwischen einer Treppenkante und dem Boden wird über den Richtungsvektor der Kante und einen Normalenvektor des Bodens berechnet. Welche Funktion steht in der Formel? \\ \kreuz{Kosinus} \\ \kreuz{Sinus}
\teil Berechne den Winkel, unter dem die Gerade $g\colon \vec x = (1 \mid 0 \mid 2) + t \cdot (2 \mid 1 \mid 2)$, $t \in \mathbb{R}$, die xy-Ebene schneidet. $\varphi \approx$ \leerfeld[$^\circ$]
\teil Ein Pyramidendach hat die Traufecke $P(3 \mid 3 \mid 2)$ und die Spitze $S(0 \mid 0 \mid 5)$; die Grundfläche des Hauses liegt in der xy-Ebene. Berechne den Neigungswinkel der Dachkante PS gegen die Grundfläche. $\varphi \approx$ \leerfeld[$^\circ$]
\teil Ein Stahlseil ist von $A(1 \mid 2 \mid 0)$ nach $B(4 \mid 6 \mid 5)$ gespannt, die xy-Ebene ist die Horizontale. Berechne den Neigungswinkel des Seils gegen die Horizontale über den Winkel zwischen $\overrightarrow{AB}$ und seiner Projektion in die xy-Ebene. $\varphi \approx$ \leerfeld[$^\circ$]
\teil Das Tragseil einer Hängebrücke wird in einem Querschnitt durch einen Kreisbogen durch $A(-30 \mid 13)$, $B(30 \mid 13)$ und den tiefsten Punkt $C(0 \mid 3)$ beschrieben; der Kreis hat den Mittelpunkt $M(0 \mid 53)$, 1 LE entspricht 1 m. Berechne die Länge des Seils zwischen A und B. \hfill (Abitur 2023 LK)
\end{teile}
\end{aufgabe}

%% TODO Titel: Ich-kann-Satz fehlt in der Bank (Platzhalter: Kettenname) – Nr. 17
%% TODO Anweisung: ein Satz über der ersten Teilaufgabe fehlt in der Bank – Nr. 17
\begin{aufgabe}{Winkel gegen Ebenen und Bogenmaße}
\begin{teile}
\teil Eine gerade Pyramide hat ihre quadratische Grundfläche in der xy-Ebene, eine Ecke $F(3 \mid 3 \mid 0)$ und die Spitze $S(0 \mid 0 \mid 6)$. Sonnenlicht fällt in Richtung $(1 \mid 1 \mid -1)$. Begründe über den Vergleich der Neigungswinkel von Kante FS und Licht gegen den Boden, ob der Schatten der Spitze neben der Pyramide auf den Boden fällt.
\end{teile}
\end{aufgabe}

%% TODO Titel: Ich-kann-Satz fehlt in der Bank (Platzhalter: Kettenname) – Nr. 18
%% TODO Anweisung: ein Satz über der ersten Teilaufgabe fehlt in der Bank – Nr. 18
\begin{aufgabe}{Winkel gegen Ebenen und Bogenmaße – Fehler finden, Begründen, Anwendung}
\begin{teile}
\teil Jonas berechnet den Winkel zwischen der Geraden mit dem Richtungsvektor $(4 \mid 1 \mid 8)$ und der xy-Ebene: $\mathrm{cos}\,\varphi = \frac{8}{9}$, $\varphi \approx 27{,}3^\circ$. Finde den Fehler.
\teil Begründe, warum für den Winkel zwischen einer Geraden und einer Ebene der Sinus statt des Kosinus in der Formel steht.
\teil Ein Flugzeug fliegt im Landeanflug geradlinig in Richtung $(-20 \mid 0 \mid -1)$, die Landebahn liegt in der xy-Ebene. Ein Anflug gilt als normal, wenn der Gleitwinkel gegen die Landebahn zwischen $2{,}5^\circ$ und $3{,}5^\circ$ liegt. Ist dieser Anflug normal?
\end{teile}
\end{aufgabe}
